% \iffalse meta-comment
%
% hit-sty.dtx 是示例文档自带的宏包，装用户自留地里的东西
%
% \fi
%
% \section{示例宏包 \file{hithesis.sty}}
%
% \changes{v3.2a}{2026/09/13}{示例宏包只装示例自己要的：\pkg{xeCJKfntef}、\pkg{xcolor}、
%   \pkg{tikz}、\pkg{tabularray} 归类，\pkg{multirow}、\pkg{placeins}、\pkg{lipsum}
%   归这里；\pkg{bm} 只在 8 位数学字体下装（\pkg{unicode-math} 下 \cs{bm} 改接
%   \cs{symbf}）；\pkg{rsfs} 花体改成可缩放声明（\file{ursfs.fd} 没有 10.53937，
%   取最近的 10.95 顶上大百分之四）；\pkg{algorithm2e} 按 \cs{g\_\_hit\_chapter\_bool}
%   分支，\cs{algorithmcfname} 按语言取；示例词条改用 \option{abbr} 族；\cs{tikzstyle}
%   改写成 \cs{tikzset}，去掉 23 个落在 \env{tikzpicture} 里的分号}
%
% v3.1e 的示例目录下有个 \file{hithesis.sty}，开头写着「此文件声明不在规范中要求的
% 格式所使用的宏包。（所以，格式基本上是自由发挥的。）」。里面是术语词条、染色体
% 示意图的 \pkg{tikz} 宏、九个数学记号速记、\cs{num} 的分组空格，都是示例文档自己
% 要用的东西，用户复制一份改成自己的。
%
% 分界线是「类自己调不调」。\pkg{siunitx}、\pkg{bm}、\pkg{mathrsfs}、\pkg{rotating}、
% \pkg{listings}、\pkg{algorithm2e} 这六个，类里一句都没用到，只是替文档预先装上并
% 配好，属于自留地。原文件里的 \pkg{xeCJKfntef}、\pkg{xcolor}、\pkg{tikz} 不在这里，
% 那三个类自己要用（下划线、封面颜色、封面图形），留在类里。
%
% \pkg{bm} 与 \pkg{mathrsfs} 看着像重复，其实不是。类只提供 \cs{boldsymbol}，
% 那不是 \cs{bm} 的全等替代；\cs{mathscr} 类里虽然有，但不装 \pkg{mathrsfs} 时它
% 跟 \cs{mathcal} 出自同一套字形，两个记号会长得一模一样，v3.1e 特意用 \pkg{rsfs}
% 就是为了把它们分开。
%
% 用法是在导言区写 \verb|\usepackage{hithesis}|。文件在示例目录下，也装进了 TEXMF；
% 当前目录优先，把示例目录那份拷进自己的项目再改就行。
%
%    \begin{macrocode}
%<*hit-sty>
\NeedsTeXFormat{LaTeX2e}
\ProvidesPackage{hithesis}
  [2026/01/01 v3.2a hithesis example package]
%    \end{macrocode}
%
% \subsection{填充文字}
%
% 深圳博士中期那份示例正文拿 \cs{lipsum} 填充。这是示例自己的需要，原先由
% \file{hit-report.cls} 在「深圳博士中期」这个组合下加载，别的组合下同一个文件就
% 报 \texttt{Undefined control sequence}。填充文字的宏包归示例宏包管。
%    \begin{macrocode}
\RequirePackage{lipsum}
%    \end{macrocode}
%
% \subsection{数字与数学}
%
% 规范对数字书写的规定（4 位及以上每 3 位空 1/4 汉字）由 \pkg{siunitx} 落实。
% 除这个空格以外任何空格都是错的，包括 \verb|\:| \verb|\;| \verb|\ |。
%    \begin{macrocode}
\RequirePackage{siunitx}
% 分组间隔按《1-8 页》范例实测:正文号数字组间 3.27bp,\hspace{0.25em} 排出
% 2.77(Word 的间隔两侧各带字符间距加宽,\hspace 是裸胶),校到 0.29em。
\sisetup{group-minimum-digits=4, group-separator= \hspace{0.29em}}
\sisetup{detect-weight,detect-mode,detect-family}
%    \end{macrocode}
%
% \pkg{bm} 处理数学公式中的黑斜体。它只在 8 位数学字体下能用：\pkg{unicode-math}
% 的符号是扩展 mathchar，\pkg{bm} 处理不了，会报 \texttt{Extended mathchar used as
% mathchar}，这时把 \cs{bm} 接到 \pkg{unicode-math} 自己的 \cs{symbf} 上。
% \cs{hit@if@option@TF} 由文档类提供，读的是类选项对应的标志位。
%    \begin{macrocode}
\hit@if@option@TF{unicodemath}
  {\let\bm\symbf}
  {\RequirePackage{bm}}
%    \end{macrocode}
%
% \pkg{mathrsfs} 提供不同于 \cs{mathcal}、\cs{mathfrak} 的英文花体。
%
% 它的字形声明（\file{ursfs.fd}）只列了 5、7、8、9、10、10.95 等几档固定字号，
% 模板要的 10.53937 不在其中，\hologo{LaTeX} 取最近的 10.95 顶上，花体因此比周围大
% 百分之四，日志里每份文档报两条 size substituted。下面按尺寸区间重新声明，
% 让它从 Type1 字形缩放到要的大小。三个 \file{rsfs*.pfb} 是 \hologo{TeX}~Live 自带的。
%    \begin{macrocode}
\RequirePackage{mathrsfs}
\DeclareFontFamily{U}{rsfs}{\skewchar\font127 }
\DeclareFontShape{U}{rsfs}{m}{n}{<-6>rsfs5 <6-8>rsfs7 <8->rsfs10}{}
%    \end{macrocode}
%
% \begin{macro}{\theVector,\theMatrix,\theSet,\theDirected,\theUndirected,
%   \theNetwork,\theNode,\theDirectedEdge,\theUndirectedEdge}
% 格式和内容分离是 \hologo{LaTeX} 的长处，常用记号在这里定一次，正文里只写记号名。
% 想换全文表示网络的符号样式，改这一行就行，不用去正文里挨个找。
%
% \cs{theNetwork} 走 \pkg{rsfs} 花体，那套字形只有 A 到 Z，小写字母和数字没有对应
% 字形，写进去会被悄悄丢掉。
%    \begin{macrocode}
\newcommand{\theVector}[1]{\bm{#1}}
\newcommand{\theMatrix}[1]{\mathbb{#1}}
\newcommand{\theSet}[1]{\mathcal{#1}}
\newcommand{\theDirected}[1]{{\overrightarrow{#1}}}
\newcommand{\theUndirected}[1]{{\overline{#1}}}
\newcommand{\theNetwork}[1]{\mathscr{#1}}
\newcommand{\theNode}[1]{{\text{#1}}}
\newcommand{\theDirectedEdge}[2]{{\overrightarrow{{#1}{#2}}}}
\newcommand{\theUndirectedEdge}[2]{{\overline{{#1}{#2}}}}
%    \end{macrocode}
% \end{macro}
%
% \subsection{排版正文常用的几个包}
%
% 这两个原先由文档类装。按「类自己调不调」这条线，它们归这边：类的表格不合并
% 单元格，也不调 \cs{FloatBarrier}。
%
% \pkg{xeCJKfntef} 一度也挪到了这里，是判断错误：学位论文封面的
% \env{CJKfilltwosides} 就是它提供的，那个包由
% \file{src/pages/hit-final-cover.dtx} 自己加载。
%
% \cs{multirow} 合并表格的纵向单元格。
%    \begin{macrocode}
\RequirePackage{multirow}
%    \end{macrocode}
%
% \cs{FloatBarrier} 把还没排的浮动体就地排完，用来拦住图表跑进下一节。
% \option{below} 允许上一节的浮动体落在下一节开头。这个包不传 \option{section}
% 就不会自动在每节插栏栅，要拦得自己写 \cs{FloatBarrier}。
%    \begin{macrocode}
\RequirePackage[below]{placeins}
%    \end{macrocode}
%
% \subsection{图表与算法}
%    \begin{macrocode}
\RequirePackage{rotating}
%    \end{macrocode}
%
% 推荐的表格写法是 \pkg{tabularray}，类已经装好（连同它的 \pkg{booktabs}
% 库），直接用 \env{booktabs}、\env{longtblr} 这些环境即可，
% 见 \file{src/engine/hit-table.dtx}。
%
% 算法的宏包，注意宏包兼容性，先后顺序为 float、hyperref、algorithm2e，否则无法
% 生成算法列表。前两个由文档类加载，所以本文件必须在 \cs{documentclass} 之后加载。
% 我工算法混乱的问题详见手册，各个实验室的具体设置方法详见手册或示例中给出的地址。
%
% 下面整段按章编号，只在终稿装。开题与中期报告不排章，
% \texttt{algochapter} 会让算法按一个恒为零的章号编成「0-1」，
% 所以那边整段跳过，要用算法自己 \cs{usepackage}\{algorithm2e\}。
%
% \emph{判据不能问 \cs{chapter} 在不在。} 报告原先是 \cls{article} 底子，
% 「有没有 \cs{chapter}」恰好等于「是不是终稿」，所以那样写是对的。
% 底类换成 \cls{ctexbook} 之后 \cs{chapter} 两边都有，这个判据会翻真，报告
% 会连带装上 \pkg{algorithm2e} 与它的三个依赖，并按恒零的章号编号。改问
% 一个自家的标志位：它说的是「这一档排不排章」，与底类给不给 \cs{chapter}
% 是两件事。
%    \begin{macrocode}
\newif\ifhit@sty@chapter
\ExplSyntaxOn
\bool_if_exist:cT { g__hit_chapter_bool }
  { \bool_if:cT { g__hit_chapter_bool } { \hit@sty@chaptertrue } }
\ExplSyntaxOff
\ifhit@sty@chapter
\RequirePackage[algoruled,linesnumbered,algochapter]{algorithm2e}
\SetAlCapSty{}
%    \end{macrocode}
%
% \begin{macro}{\foocaption}
%    \begin{macrocode}
\newcommand{\foocaption}[1]{ \def\@algocf@pre@plainruled{\hrule height1.5pt depth0pt\kern\interspacetitleruled #1 \kern\interspacealgoruled\hrule height1pt depth0pt\kern\interspacetitleruled} }
\def\@algocf@post@ruled{\kern\interspacealgoruled\hrule height1.5pt\relax}%
%    \end{macrocode}
% \end{macro}
%
% 算法列表。中文名 \cs{listalgoenname} 与 \cs{hit@caption@algorithm@zh} 由文档类提供。
%    \begin{macrocode}
\newcommand{\algoenname}{Algo.}
\newfloatlist[chapter]{algoen}{aen}{\listalgoenname}{\algoenname}
\newfixedcaption{\algoencaption}{algoen}
\renewcommand{\thealgoen}{\thechapter-\arabic{algocf}}
\renewcommand{\@cftmakeaentitle}{\chapter*{\listalgoenname\@mkboth{\listalgoenname}{\listalgoenname}}
}
\hit@if@option@TF{en}
  {\renewcommand{\algorithmcfname}{\hit@term@caption@algorithm@en}}
  {\renewcommand{\algorithmcfname}{\hit@term@caption@algorithm@zh}}
\setlength\AlCapSkip{1.2ex}
\SetAlgoSkip{1pt}
\renewcommand{\algocf@captiontext}[2]{\wuhao#1\algocf@typo~\AlCapFnt{}#2}
\expandafter\ifx\csname algocf@within\endcsname\relax
\renewcommand\thealgocf{\@arabic\c@algocf}
\else
\renewcommand\thealgocf{\csname the\algocf@within\endcsname-\@arabic\c@algocf}
\fi
%    \end{macrocode}
%
% 中英文双标题一定多于一行，因此去掉单行时的判断，并将 \cs{parbox} 中标题设置为居中。
%    \begin{macrocode}
\renewcommand{\algocf@makecaption}[2]{%
  \addtolength{\hsize}{\algomargin}%
  \sbox\@tempboxa{\algocf@captiontext{#1}{#2}}%
    \hskip .5\algomargin%
    \parbox[t]{\hsize}{\centering\algocf@captiontext{#1}{#2}}%
  \addtolength{\hsize}{-\algomargin}%
}
%    \end{macrocode}
%
% \begin{macro}{\AlgoBiCaption}
% 直接取出自定义的中英文标题条目加入到真正的 \cs{caption} 中。
%    \begin{macrocode}
\newcommand{\AlgoBiCaption}[2]{%
   \caption[#1]{\protect\setlength{\baselineskip}{1.5em}#1 \protect \\ Algo.\thealgocf~#2}
   \addcontentsline{aen}{algoen}{\protect\numberline{\thealgoen}{#2}}
}
\else
\PackageWarningNoLine{hithesis}
  {这一档不排章，跳过 algorithm2e 那一段（它按章编号）。\MessageBreak
   要排算法请自己 \protect\usepackage{algorithm2e}}
\fi
%    \end{macrocode}
% \end{macro}
%
% 排版源码所使用的环境可以有两种，\pkg{listings} 和 \pkg{minted}。用 \pkg{minted}
% 要在编译时加 \texttt{-shell-escape} 选项，并装 pygmentize 工具给代码高亮。
%    \begin{macrocode}
\RequirePackage{listings}
\lstset{
% basicstyle=\small\ttfamily,
columns=flexible,
breaklines=true
}
% \RequirePackage{minted}
%    \end{macrocode}
%
% \subsection{缩略语词条}
%
% 缩略语词条填在 \option{abbr} 族里，键名就是短形，值收 \option{zh} 与
% \option{en}；只有一种语言时写裸值，含汉字归中文、不含归英文。索引形式由
% \option{zh} 与 \option{en} 直接来，中文那份的拼音排序键自动查表推。
%
% v3.1e 是每个中文术语手搓一个宏（\cs{gtssbp}、\cs{gscna}、\cs{gscnas}），
% 把 \cs{gls} 与两句 \cs{sindex} 打包在一起。作者得记住用那个宏而不是普通命令，
% 写错就漏索引，单复数还要各写一个。现在正文一律写 \cs{abbr} 与 \cs{abbrs}。
%
% 这一段两档都能编：\option{abbr} 族只有终稿有，所以先看 \cs{hitsetup}
% 认不认它。
%    \begin{macrocode}
\ExplSyntaxOn
\keys_if_exist:nnT { hit } { abbr }
  {
    \hitsetup
      {
        abbr =
          {
            TSSBP = { zh = 树结构折筷过程 ,
                      en = Tree-structured~Stick-breaking~process } ,
            SCNA  = { zh = 体细胞拷贝数变异 ,
                      en = Somatic~copy~number~alteration } ,
            TSE   = { Taylor~Series~Expansion } ,
            SVM   = { Support~Vector~Machine } ,
            ML    = { Machine~Learning } ,
            CO    = { Convex~Optimization } ,
          }
      }
  }
\ExplSyntaxOff
%    \end{macrocode}
%
% \subsection{染色体示意图}
%
% \file{body/introduction.tex} 里有一张染色体示意图，下面这批样式和宏是画它用的。
% 画的是固定内容，\cs{rcell} 直接画 A~C~T~C 四个碱基。留在这里当画图的例子看。
%
% \pkg{tikz} 与 \pkg{xcolor} 由文档类加载。原文件用的 \cs{tikzstyle} 已经废弃，
% 这里改成 \cs{tikzset}。
%
% 四个颜色只给下面这批 \cs{tikzset} 用。
%    \begin{macrocode}
\definecolor{colorzero}{rgb}{0, 0, 0}
\definecolor{colorone}{rgb}{1, 0, 0}
\definecolor{colortwo}{rgb}{0, 0, 1}
\definecolor{colorthree}{rgb}{0, 1, 0}
%    \end{macrocode}
%
% 此处可以定义一些 \pkg{tikz} 全局样式，例如
% \verb|\tikzset{nodestyle/.style={circle, fill=gray!60}}|、
% \verb|\tikzset{edgestyle/.style={-latex}}|。
%    \begin{macrocode}
% \pkg{tikz} 也只有终稿加载，报告没有，所以这批样式同样要看在不在。
% 下面 \cs{gseg} 那批宏本身只是定义，缺 \pkg{tikz} 也定义得出来，用到才会炸。
\ifdefined\tikzset
\tikzset{maternal/.style={colorone}}
\tikzset{paternal/.style={colortwo}}
\tikzset{variant/.style={colorthree!80!colorzero}}
\tikzset{reference/.style={colorzero}}
\tikzset{aallele/.style={colorzero,rotate=90}}
\tikzset{ballele/.style={colorthree!80!colorzero,rotate=90}}
\tikzset{refseg/.style={colorzero,draw=colorzero, opacity=0.2}}
\tikzset{mseg/.style={colorone,draw=colorone, opacity=0.2}}
\tikzset{pseg/.style={colortwo,draw=colortwo, opacity=0.2}}
\tikzset{vseg/.style={colorthree!80!colorzero,draw=colorthree!80!colorzero, opacity=0.6}}
\tikzset{bncell/.style={draw=colorzero,opacity=0.2,line width=2pt, rounded corners=1pt}}
\tikzset{btcell/.style={draw=colorone,opacity=0.6, line width=2pt, rounded corners=1pt}}
\tikzset{tncell/.style={colorzero,opacity=0.9}}
\tikzset{ttcell/.style={colorone,opacity=0.6}}
\tikzset{tscell/.style={colorzero}}
\tikzset{refcell/.style={colorzero}}
\tikzset{evolve/.style={->,draw=colortwo,opacity=0.3,line width=1.5pt}}
\tikzset{fakeevolve/.style={->,draw=colorzero,opacity=0.3,line width=1.5pt}}
\tikzset{refline/.style={dashed,draw=colorzero,line width=1pt}}
\tikzset{tnline/.style={dashed,draw=colorzero,opacity=0.3,line width=1pt}}
\fi
%    \end{macrocode}
%
% \begin{macro}{\gseg,\gsegr}
% 一段染色体，带两端的延长线。\cs{gsegr} 是右偏移的版本。
%
% 下面各个 cell 宏里调用 \cs{gseg} 之后的那个分号，v3.1e 原文是有的。它落在
% \env{tikzpicture} 里没有路径可收，每份文档报 23 条 \texttt{Missing character:
% There is no ; in font nullfont}，排出来什么也没有。这里去掉。
%    \begin{macrocode}
\newcommand{\gseg}[9]{%
  \pgfmathsetmacro{\sstartx}{#1}
  \pgfmathsetmacro{\slengx}{#2}
  \pgfmathsetmacro{\sy}{#3}
  \pgfmathsetmacro{\sdy}{#4}
  \pgfmathsetmacro{\sdx}{#5}
  \pgfmathsetmacro{\sdxh}{#7}
  \pgfmathsetmacro{\sdxt}{#8}
  \fill[#6](\sstartx,\sy)--(\sstartx-\sdx,\sy+\sdy)--
  (\slengx+\sstartx+1.5-\sdx,\sy+\sdy)--(\slengx+\sstartx+1.5,\sy)--
  (\slengx+\sstartx+1.5-\sdx,\sy-\sdy)--(\sstartx-\sdx,\sy-\sdy)--cycle;
  \draw[#9] (\sstartx-\sdxh,\sy) -- (\sstartx, \sy);
  \draw[#9] (\slengx+\sstartx+1.5, \sy) -- (\slengx+\sstartx+1.5+\sdxt,\sy);
}
\newcommand{\gsegr}[9]{%
  \pgfmathsetmacro{\sstartx}{#1}
  \pgfmathsetmacro{\slengx}{#2}
  \pgfmathsetmacro{\sy}{#3}
  \pgfmathsetmacro{\sdy}{#4}
  \pgfmathsetmacro{\sdx}{#5}
  \pgfmathsetmacro{\sdxh}{#7}
  \pgfmathsetmacro{\sdxt}{#8}
  \fill[#6](\sstartx-0.5,\sy)--(\sstartx+\sdx-0.5,\sy+\sdy)--
  (\slengx+\sstartx+1.5+\sdx-0.5,\sy+\sdy)--(\slengx+\sstartx+1.5-0.5,\sy)--
  (\slengx+\sstartx+1.5+\sdx-0.5,\sy-\sdy)--(\sstartx+\sdx-0.5,\sy-\sdy)--cycle;
  \draw[#9] (\sstartx-\sdxh-0.5,\sy) -- (\sstartx-0.5, \sy);
  \draw[#9] (\slengx+\sstartx+1.5-0.5, \sy) -- (\slengx+\sstartx+1.5+\sdxt-0.5,\sy);
}
%    \end{macrocode}
% \end{macro}
%
% \begin{macro}{\rcell,\ncell}
% 参考序列与正常细胞。
%    \begin{macrocode}
\newcommand{\rcell}[2]{%
  \pgfmathsetmacro{\x}{#1}
  \pgfmathsetmacro{\y}{#2}
  %\node at (\x+10, \y) {Reference};
  \draw (\x+1,\y) node[aallele]{A};
  \draw (\x+2,\y) node[aallele]{C};
  \draw (\x+3,\y) node[aallele]{T};
  \draw (\x+4,\y) node[aallele]{C};
  \gseg{\x}{4}{\y}{0.2}{0.5}{refseg}{1.5}{1}{reference}
}
\newcommand{\ncell}[2]{%
  \pgfmathsetmacro{\x}{#1}
  \pgfmathsetmacro{\y}{#2}
  %\node [maternal] at (\x+8, \y) {M};
  %\node [paternal] at (\x+8, \y-0.5) {P};
  \draw[bncell](\x-2,\y+0.5)--(\x+7,\y+0.5)--
  (\x+7,\y-1)--(\x-2,\y-1)--cycle;
  \draw (\x+1,\y) node[aallele]{A};
  \draw (\x+2,\y) node[ballele]{G};
  \draw (\x+3,\y) node[aallele]{T};
  \draw (\x+4,\y) node[aallele]{C};
  \gseg{\x}{4}{\y}{0.2}{0.5}{mseg}{1.5}{1}{maternal}
  \draw (\x+1,\y-0.5) node[ballele]{T};
  \draw (\x+2,\y-0.5) node[aallele]{C};
  \draw (\x+3,\y-0.5) node[aallele]{T};
  \draw (\x+4,\y-0.5) node[ballele]{A};
  \gseg{\x}{4}{\y-0.5}{0.2}{0.5}{pseg}{1.5}{1}{paternal}
}
%    \end{macrocode}
% \end{macro}
%
% \begin{macro}{\tcellone,\tcelltwo,\tcellthree,\tcellfour,\tcellfive}
% 五种肿瘤细胞。
%    \begin{macrocode}
\newcommand{\tcellone}[2]{%
  \pgfmathsetmacro{\x}{#1}
  \pgfmathsetmacro{\y}{#2}
  %\node [maternal] at (\x+8, \y) {M};
  %\node [maternal] at (\x+8, \y-0.5) {M};
  %\node [paternal] at (\x+8, \y-1) {P};
  \draw[btcell](\x-2,\y+0.5)--(\x+7,\y+0.5)--
  (\x+7,\y-1.5)--(\x-2,\y-1.5)--cycle;
  \draw (\x+1,\y) node[aallele]{A};
  \draw (\x+2,\y) node[ballele]{G};
  \draw (\x+3,\y) node[aallele]{T};
  \draw (\x+4,\y) node[aallele]{C};
  \gseg{\x}{4}{\y}{0.2}{0.5}{mseg}{1.5}{1}{maternal}
  \draw (\x+1,\y-0.5) node[aallele]{A};
  \draw (\x+2,\y-0.5) node[ballele]{G};
  \draw (\x+3,\y-0.5) node[aallele]{T};
  \draw (\x+4,\y-0.5) node[aallele]{C};
  \gseg{\x}{4}{\y-0.5}{0.2}{0.5}{mseg}{1.5}{1}{maternal}
  \draw (\x+1,\y-1) node[ballele]{T};
  \draw (\x+2,\y-1) node[aallele]{C};
  \draw (\x+3,\y-1) node[aallele]{T};
  \draw (\x+4,\y-1) node[ballele]{A};
  \gseg{\x}{4}{\y-1}{0.2}{0.5}{pseg}{1.5}{1}{paternal}
}
\newcommand{\tcelltwo}[2]{%
  \pgfmathsetmacro{\x}{#1}
  \pgfmathsetmacro{\y}{#2}
  %\node [maternal] at (\x+8, \y) {M};
  %\node [maternal] at (\x+8, \y-0.5) {M};
  %\node [paternal] at (\x+8, \y-1) {P};
  \draw[btcell](\x-2,\y+0.5)--(\x+7,\y+0.5)--
  (\x+7,\y-1.5)--(\x-2,\y-1.5)--cycle;
  \draw (\x+1,\y)  node[aallele]{A};
  \draw (\x+2,\y)  node[ballele]{G};
  \draw (\x+3,\y)  node[aallele]{T};
  \draw (\x+4,\y)  node[aallele]{C};
  \gseg{\x}{4}{\y}{0.2}{0.5}{mseg}{1.5}{1}{maternal}
  \draw (\x+1,\y-0.5) node[aallele]{A};
  \draw (\x+2,\y-0.5) node[ballele]{G};
  \draw (\x+3,\y-0.5) node[aallele]{T};
  \draw (\x+4,\y-0.5) node[ballele]{G};
  \gseg{\x}{4}{\y-0.5}{0.2}{0.5}{mseg}{1.5}{1}{maternal}
  \draw (\x+1,\y-1) node[ballele]{T};
  \draw (\x+2,\y-1) node[aallele]{C};
  \draw (\x+3,\y-1) node[aallele]{T};
  \draw (\x+4,\y-1) node[ballele]{A};
  \gseg{\x}{4}{\y-1}{0.2}{0.5}{pseg}{1.5}{1}{paternal}
}
\newcommand{\tcellthree}[2]{%
  \pgfmathsetmacro{\x}{#1}
  \pgfmathsetmacro{\y}{#2}
  %\node [maternal] at (\x+12, \y) {M};
  %\node [paternal] at (\x+12, \y-0.5) {P};
  \draw[btcell](\x-2,\y+0.5)--(\x+11,\y+0.5)--
  (\x+11,\y-1)--(\x-2,\y-1)--cycle;
  \draw (\x+1,\y) node[aallele]{A};
  \draw (\x+2,\y) node[ballele]{G};
  \gseg{\x}{2}{\y}{0.2}{0.5}{mseg}{1.5}{0}{maternal}
  \gseg{\x+4}{0}{\y}{0.2}{0.5}{vseg}{0.5}{0.5}{variant}
  \draw (\x+7,\y) node[aallele]{T};
  \draw (\x+8,\y) node[aallele]{C};
  \gseg{\x+6}{2}{\y}{0.2}{0.5}{mseg}{0}{1}{maternal}
  \draw (\x+1,\y-0.5) node[ballele]{T};
  \draw (\x+2,\y-0.5) node[aallele]{C};
  \draw (\x+3,\y-0.5) node[aallele]{T};
  \draw (\x+4,\y-0.5) node[ballele]{A};
  \gseg{\x}{4}{\y-0.5}{0.2}{0.5}{pseg}{1.5}{1}{paternal}
}
\newcommand{\tcellfour}[2]{%
  \pgfmathsetmacro{\x}{#1}
  \pgfmathsetmacro{\y}{#2}
  %\node [maternal] at (\x+18, \y) {M};
  %\node [paternal] at (\x+18, \y-0.5) {P};
  \draw[btcell](\x-2,\y+0.5)--(\x+15,\y+0.5)--
  (\x+15,\y-1)--(\x-2,\y-1)--cycle;
  \draw (\x+1,\y) node[aallele]{A};
  \draw (\x+2,\y) node[ballele]{G};
  \gseg{\x}{2}{\y}{0.2}{0.5}{mseg}{1.5}{0}{maternal}
  \gseg{\x+4}{0}{\y}{0.2}{0.5}{vseg}{0.5}{0.5}{variant}
  \draw (\x+7,\y) node[aallele]{T};
  \gseg{\x+6}{1}{\y}{0.2}{0.5}{mseg}{0}{0}{maternal}
  \gseg{\x+9}{0}{\y}{0.2}{0.5}{vseg}{0.5}{0.5}{variant}
  \draw (\x+12,\y) node[aallele]{C};
  \gseg{\x+11}{1}{\y}{0.2}{0.5}{mseg}{0}{1}{maternal}
  \draw (\x+1,\y-0.5) node[ballele]{T};
  \draw (\x+2,\y-0.5) node[aallele]{C};
  \draw (\x+3,\y-0.5) node[aallele]{T};
  \draw (\x+4,\y-0.5) node[ballele]{A};
  \gseg{\x}{4}{\y-0.5}{0.2}{0.5}{pseg}{1.5}{1}{paternal}
}
\newcommand{\tcellfive}[2]{%
  \pgfmathsetmacro{\x}{#1}
  \pgfmathsetmacro{\y}{#2}
  %\node [maternal] at (\x+8, \y) {M};
  %\node [maternal] at (\x+8, \y-0.5) {M};
  %\node [paternal] at (\x+8, \y-1) {P};
  \draw[btcell](\x-2,\y+0.5)--(\x+9.5,\y+0.5)--
  (\x+9.5,\y-1.5)--(\x-2,\y-1.5)--cycle;
  \draw (\x+1,\y) node[aallele]{A};
  \draw (\x+2,\y) node[ballele]{G};
  \draw (\x+3,\y) node[aallele]{T};
  \draw (\x+4,\y) node[aallele]{C};
  \gseg{\x}{4}{\y}{0.2}{0.5}{mseg}{1.5}{1}{maternal}
  \draw (\x+1,\y-0.5) node[aallele]{A};
  \draw (\x+2,\y-0.5) node[ballele]{G};
  \draw (\x+3,\y-0.5) node[aallele]{T};
  \draw (\x+4,\y-0.5) node[aallele]{C};
  \gseg{\x}{4}{\y-0.5}{0.2}{0.5}{mseg}{1.5}{1}{maternal}
  \draw (\x+1,\y-1) node[ballele]{T};
  \gseg{\x}{1}{\y-1}{0.2}{0.5}{pseg}{1.5}{0}{paternal}
  \draw (\x+4.5,\y-1) node[ballele]{A};
  \draw (\x+5.5,\y-1) node[aallele]{T};
  \draw (\x+6.5,\y-1) node[aallele]{C};
  \gsegr{\x+3.5}{3}{\y-1}{0.2}{0.5}{pseg}{0.5}{1.5}{paternal}
}
%    \end{macrocode}
% \end{macro}
%
% \subsection{正文里排控制序列}
%
% \verb|\cs{foo}| 排出 \verb|\foo|，示例正文讲命令用法时在用。
%    \begin{macrocode}
\def\cmd#1{\cs{\expandafter\cmd@to@cs\string#1}}
\def\cmd@to@cs#1#2{\char\number`#2\relax}
\DeclareRobustCommand\cs[1]{\texttt{\char`\\#1}}
%</hit-sty>
%    \end{macrocode}
%
% \Finale
